% titel: Blatt 1, A1: Messgenauigkeit
% Teil II: der durchgerechnete Weg.

\subsection*{Gegeben und gesucht}
Fünf Wägungen (in g): $26.518,\ 26.526,\ 26.504,\ 26.511,\ 26.534$. Gesucht
Mittelwert $\bar x$ und Standardabweichung $s$.

\subsection*{Rechnung}
Mittelwert:
\begin{equation}
  \bar x = \frac{26.518+26.526+26.504+26.511+26.534}{5}
  = \frac{132.593}{5} = 26.5186\ \mathrm{g}.
\end{equation}
Abweichungen quadrieren:
\begin{align*}
  (0.0006)^2 &= 3.6\cdot10^{-7}, &
  (0.0074)^2 &= 5.476\cdot10^{-5}, \\
  (-0.0146)^2 &= 2.1316\cdot10^{-4}, &
  (-0.0076)^2 &= 5.776\cdot10^{-5}, \\
  (0.0154)^2 &= 2.3716\cdot10^{-4}.
\end{align*}
Quadratsumme: $\Sigma(x_i-\bar x)^2 = 5.62\cdot10^{-4}\ \mathrm{g^2}$.
\begin{equation}
  s = \sqrt{\frac{5.62\cdot10^{-4}}{5-1}} = \sqrt{1.405\cdot10^{-4}}
  \approx 0.01185\ \mathrm{g} \approx 0.012\ \mathrm{g}.
\end{equation}

\begin{ergebnis}
$\bar x = 26.519\ \mathrm{g},\qquad s \approx 0.012\ \mathrm{g}.$
\end{ergebnis}

\begin{kontrolle}
Die Werte streuen $\pm 0.015$ g um 26.52 g; eine Standardabweichung von
rund 0.012 g stimmt mit dieser sichtbaren Streuung überein. $s$ ist deutlich
kleiner als die Messwerte selbst --- die Wägung ist präzise.
\end{kontrolle}

\begin{merksatz}
Mehrfachmessung: erst Mittelwert, dann die Abweichungen quadriert mitteln
($n-1$) und die Wurzel ziehen. So beschreibt man eine Messreihe mit einem
Wert und einer Unsicherheit.
\end{merksatz}
